\chapter{Derivations for Bayes Nets as Models of Rationality}
\label{annex:p16-rationality-4-2}

\begin{annexitem}{annex:p16.rat42-product}{thm:p16-rat42-product}{Normalized tables, PMC, and parameter economy}
\begin{proof}
Choose a topological ordering of the finite DAG, whose existence is
\cref{thm:p16-causal-topological-order}; it is a graph property and
requires no causal interpretation. The product of nonnegative tables
is nonnegative. Sum first over the last node in this order. Its variable
occurs in its own factor but no remaining child's factor, and its table
sums to one. Repeat backwards to obtain total mass one.
The same summation shows that an ancestral marginal retains exactly
its ancestral factors: vertices outside the ancestral set can be
removed in reverse topological order without removing a retained parent.

For node \(A_i\), let \(R\) consist of the node and every vertex that
is not its strict descendant. The set \(R\) is ancestral. Indeed,
an ancestor outside \(R\) would be a descendant of \(A_i\), and a
path from it to a member of \(R\) would make that member a descendant
too, or create a cycle if it were \(A_i\). Within \(R\), \(A_i\)
has no children. Consequently
\[
 P_R(a_i,a_{R\setminus\{i\}})
   =k_i(a_i\mid a_{pa(A_i)})
       \prod_{j\in R\setminus\{i\}}k_j(a_j\mid a_{pa(A_j)}),
\]
and the second product is the marginal on the remaining vertices.
On every positive remaining assignment the conditional of \(A_i\)
therefore equals its local table. Summing over nondescendant nonparents
also gives that table as its conditional on its parents, whenever the
parent assignment has positive probability. The displayed factorization
is precisely joint independence from the nondescendant nonparents
given the parents. This proves PMC, including laws with zeros.

Conversely, every predecessor of a node in a topological order is a
nondescendant. PMC and decomposition make the conditional of that
node on the positive predecessor assignment equal to its conditional
on its parents. The ordinary chain rule for a positive prefix thus
multiplies the required local conditional. Induct on prefix length:
if the shorter prefix has zero probability, its already established
product is zero and every longer prefix also has mass zero; otherwise
multiply by the ordinary conditional. A parent assignment inside a
positive prefix is itself positive. Arbitrary normalized tables chosen
at null parent assignments therefore do not affect this induction.
The full prefix is the required factorization.

A binary joint law has \(2^n\) entries and one normalization constraint,
so \(2^n-1\) free parameters in its unrestricted simplex. For a node
with \(k\) binary parents, each of its \(2^k\) parent configurations
needs one parameter for the probability of its positive value; its
negative probability is the complement. Summing gives
\(\sum_i2^{|pa(A_i)|}\). A chain has one root parameter and two
per later node; a fully connected topological order has
\(1+2+\cdots+2^{n-1}=2^n-1\), showing the equality case.
The assertion counts unrestricted table specifications. If parent
strata are null, some entries are unidentifiable from the joint law;
the count does not claim otherwise. Marginalization and positive-event
conditioning now follow directly by summing this normalized mass function.
\end{proof}
\end{annexitem}

\begin{annexitem}{annex:p16.rat42-semigraphoid}{thm:p16-rat42-semigraphoid}{The four finite independence operations}
\begin{proof}
Fix a conditioning tuple \(c\) of positive probability. Symmetry follows
from \(P(a,b\mid c)=P(a\mid c)P(b\mid c)\), since multiplication
commutes. For decomposition, sum
\(P(a,b,d\mid c)=P(a\mid c)P(b,d\mid c)\) over \(d\).
For weak union, when \(P(d\mid c)>0\), divide the same identity
by \(P(d\mid c)\) to obtain
\[
 P(a,b\mid d,c)=P(a\mid c)P(b\mid d,c).
\]
Summing over \(b\) gives \(P(a\mid d,c)=P(a\mid c)\), so the
right side is the required product of conditional marginals.

For contraction and a positive \((b,c)\) event, the two assumptions give
\[
\begin{aligned}
 P(a,b,d\mid c)
 &=P(b\mid c)P(a,d\mid b,c)\\
 &=P(b\mid c)P(a\mid b,c)P(d\mid b,c)\\
 &=P(a\mid c)P(b,d\mid c).
\end{aligned}
\]
If \(P(b\mid c)=0\), both sides are zero and no conditional on that
null event is needed. This proves the joint independence asserted
by contraction on every positive conditioning stratum.
\end{proof}
\end{annexitem}

\begin{annexitem}{annex:p16.rat42-three-networks}{thm:p16-rat42-three-networks}{The original graph comparison and the collider arithmetic}
\begin{proof}
The independence identity gives
\(P(a,b,c)=P(c)P(a\mid c)P(b\mid c)\) on positive \(c\).
The zero-\(c\) cells are zero, so arbitrary normalized tables on
such rows give the same factorization. This is the fork representation.
For the chain use
\(P(b)P(c\mid b)=P(b,c)=P(c)P(b\mid c)\) on their positive
strata, with zero products elsewhere. Multiplying by
\(P(a\mid c)\) gives the identical law. Conversely either product
has \(A\perp B\mid C\) by direct cancellation on positive
\(c\) strata, or by PMC and symmetry. Thus orientation alone does
not distinguish these two representations of the same constraint.

In the collider, its factorization is
\(P(a,b,c)=P(a)P(b)P(c\mid a,b)\). In particular the two
positive-\(A\) cells in the statement are \(ab\alpha\) and
\(a\bar b\beta\). Summing yields
\[
 P(\mathrm A\mid\mathrm B,\mathrm C)
 =\frac{a\alpha}{a\alpha+\bar a\gamma},
\]
and
\[
 P(\mathrm A\mid\mathrm C)
 =\frac{a(b\alpha+\bar b\beta)}
 {a(b\alpha+\bar b\beta)+\bar a(b\gamma+\bar b\delta)},
\]
whenever the corresponding denominators are positive. At the supplied
values the first is \((9/10)/(9/10+7/10)=9/16\); the numerator
of the second is \(17/40\) and its denominator \(25/40\),
giving \(17/25\). Because \(\mathrm B,\mathrm C\) has positive
probability, their inequality disproves conditional independence.

If all four table entries equal \(2/5\), then
\(P(a,b,\mathrm C)=(2/5)P(a)P(b)\) and
\(P(a,b,\neg\mathrm C)=(3/5)P(a)P(b)\). Both conditional
parent laws are their original independent product, proving
the stated counterexample to an unconditional collider-dependence claim.
\end{proof}
\end{annexitem}

\begin{annexitem}{annex:p16.rat42-kl-update}{thm:p16-rat42-kl-update}{Full finite minimization under event and conditional constraints}
\begin{proof}
First, relative entropy is nonnegative, with equality precisely for
equal mass functions. For finite value use \(\log z\leq z-1\)
at \(z=P(\omega)/Q(\omega)\) on the positive \(Q\)-support.
Summing gives
\[
 D(Q\|P)\geq
  \sum_{Q>0}(Q(\omega)-P(\omega))
  =1-\sum_{Q>0}P(\omega)\geq0.
\]
Equality requires equality in every logarithmic inequality and no
remaining prior mass, hence \(Q=P\). Infinite-divergence cases cannot
minimize when a finite feasible candidate exists.

If \(Q(\mathrm E)=1\) and \(R=P(\cdot\mid\mathrm E)\),
every finite-divergence feasible \(Q\) is supported on the prior's
positive part of \(E\). Directly replacing \(P\) by
\(P(\mathrm E)R\) in its nonzero terms gives
\[
              D(Q\|P)=D(Q\|R)-\log P(\mathrm E).
\]
The nonnegative first term is uniquely minimized by \(Q=R\).

For the four-cell conditional constraint write
\(z=Q(\mathrm A)>0\), and let
\(w=Q(\mathrm C\mid\neg\mathrm A)\) when \(z<1\), choosing
any \(w\) at \(z=1\). Splitting the four terms by the value of
\(A\) gives exactly
\[
 D(Q\|P)=D_{\rm Ber}(z\|s)+zD_{\rm Ber}(t\|u)
                         +(1-z)D_{\rm Ber}(w\|v).
\]
Here \(D_{\rm Ber}\) is the displayed two-value relative entropy,
with the zero-term conventions at endpoints. Put
\(d=D_{\rm Ber}(t\|u)\),
\(Z=\bar s+s e^{-d}\), and \(s'=s e^{-d}/Z\).
Because the prior has all four cells positive, \(0<s'<1\)
and \(d\) is finite even for \(t=0,1\). Another exact expansion is
\[
 D_{\rm Ber}(z\|s)+zd
                     =D_{\rm Ber}(z\|s')-\log Z.
\]
Thus the full objective is minimized uniquely at \(z=s'\) and
\(w=v\). The constraint already fixes the \(A\)-conditional to
\(t\); these three numbers determine all four cells and prove
both existence and uniqueness of the claimed posterior.

For \(t=1\), \(d=\log(1/u)\) and \(e^{-d}=u\).
The new \(\mathrm A,\mathrm C\) cell is
\(su/(\bar s+su)\), the \(\mathrm A,\neg\mathrm C\) cell is
zero, and each \(\neg\mathrm A\) cell is its prior mass divided
by \(\bar s+su\). That denominator is
\(P(\neg\mathrm A\lor\mathrm C)\), so this is ordinary
conditioning on the material conditional event. For \(t<1\)
the permitted \(\mathrm A,\neg\mathrm C\) cell can instead
be positive, so the procedures differ.

The direction matters: reversing \(P,Q\) would assign infinite
divergence whenever the prior puts positive mass where the certain-evidence
posterior must put zero. This theorem does not minimize that reverse
quantity. Also \(D((.9,.1)\|(.5,.5))\) differs from
\(D((.5,.5)\|(.9,.1))\), showing lack of symmetry and hence
failure to be a metric. This mathematical minimization establishes
the consequence of a specified conservativity criterion, not the
criterion's universal normative authority or a human mechanism.
\end{proof}
\end{annexitem}

\begin{annexitem}{annex:p16.rat42-confirmation}{thm:p16-rat42-confirmation}{Bayes' rule with its endpoint and strictness conditions}
\begin{proof}
The total probability of the evidence is \(hp+\bar hq\), and
the joint probability of hypothesis and evidence is \(hp\).
Divide to obtain the posterior. If \(p>0\), divide numerator
and denominator by \(p\) to obtain the inverse-likelihood-ratio
form. Subtracting the prior gives
\[
 \frac{hp}{hp+\bar hq}-h
       =\frac{h\bar h(p-q)}{hp+\bar hq}.
\]
The denominator is positive by the evidence assumption and
\(h\bar h>0\), so the sign is exactly that of \(p-q\).
The zero likelihood endpoints follow directly from the first
fraction. If both likelihoods vanish, its denominator is zero
and ordinary conditionalization is undefined.
\end{proof}
\end{annexitem}

\begin{annexitem}{annex:p16.rat42-witness}{thm:p16-rat42-witness}{Marginalizing reliability and comparing the two complete models}
\begin{proof}
The root reliability is independent of the hypothesis. Summing
its two conditional report probabilities gives
\(p=r+a\bar r\) under \(\mathrm H\), and
\(q=a\bar r\) under \(\neg\mathrm H\). Therefore
\[
 P(\mathrm H\mid\mathrm{Rep})
  =\frac{h(r+a\bar r)}{h(r+a\bar r)+\bar h a\bar r}
  =\frac{h(r+a\bar r)}{hr+a\bar r}.
\]
The likelihood difference is \(r\), so the confirmation statement
and its irrelevant and perfectly reliable limits follow from
\cref{thm:p16-rat42-confirmation}, with its positive-evidence condition.
When \(r=0=a\), a positive report has zero probability and no
posterior is asserted.

For separate reliability roots the two reports factor given
the hypothesis after summing each independent reliability.
The two-positive-report likelihoods are consequently \(p^2,q^2\).
For shared reliability, if the source is reliable both reports
are positive with probability one under \(H\) and zero under
\(\neg H\). If unreliable, conditional independence of its
two randomizations gives probability \(a^2\) under either truth
value. Marginalizing the common reliability thus gives
\(p_D=r+a^2\bar r\), \(q_D=a^2\bar r\).

At interior \(r,a\), both negative-hypothesis likelihoods are
positive. The posterior for fixed \(0<h<1\) is strictly
increasing in the ordinary Bayes factor \(L=p/q\), as follows
by cross-multiplying \(hL/(hL+\bar h)\) at two positive \(L\)'s.
The separate minus shared factors have sign determined by
\[
 (r+a\bar r)^2-\bar r(r+a^2\bar r)
             =r\bigl[r+(2a-1)\bar r\bigr].
\]
For shared confirmation to be larger this bracket must be negative.
It is equivalent to
\(2(1-a)r<1-2a\), which is possible precisely in the stated
\(a<1/2\) range with the given strict threshold. Equality
gives equal posteriors. For \(a=1/2\), the bracket is \(r>0\),
so separate reliability is strictly better.

With \(r=a=1/10\), the one-report likelihoods are \(19/100\)
and \(9/100\). For \(h=1/2\), separate reliability gives
\(19^2/(19^2+9^2)=361/442\). The shared likelihoods are
\(109/1000\) and \(9/1000\), giving \(109/118\).
Their one-report marginals nevertheless both remain
\(hr+a\bar r\). These calculations prove the comparison
under equal specified single-report information.
\end{proof}
\end{annexitem}

\begin{annexitem}{annex:p16.rat42-validity}{thm:p16-rat42-validity}{What certainty and a specified soft update really imply}
\begin{proof}
The event \(\mathrm A\land(\neg\mathrm A\lor\mathrm C)\)
is exactly \(\mathrm A\land\mathrm C\). If it has positive
probability, conditioning on it makes \(\mathrm C\) certain.
This is fixed material modus ponens, irrespective of the prior
weights on the other assignments.

For the partial case, probability one for the material major
means \(P(\mathrm A,\neg\mathrm C)=0\). Because
\(P(\mathrm A)=s>0\), this gives
\(P(\mathrm C\mid\mathrm A)=1\). The retained conditionals
and new mixing weight \(s'\) yield
\[
 P(\mathrm C)=s+(1-s)t,
 \qquad P'(\mathrm C)=s'+(1-s')t.
\]
Their difference is \((s'-s)(1-t)\). It vanishes at \(t=1\),
is positive for \(t<1\), and the second expression is below
one precisely when \((1-s')(1-t)>0\). These are the claimed
strictness conditions, not an unrestricted rule for all posterior changes.

In denying the antecedent, the two learned material premises
have conjunction
\((\neg\mathrm A\lor\mathrm C)\land\neg\mathrm A
 =\neg\mathrm A\).
The proposed conclusion \(\neg\mathrm C\) is not entailed,
since \(\neg\mathrm A\land\mathrm C\) is a countervaluation.
In the first supplied law, total probability gives
\(P(\mathrm C)=\tfrac12(.9)+\tfrac12(.2)=.55\),
so \(P(\neg\mathrm C)=.45\); conditioning on
\(\neg\mathrm A\) gives \(.8\). In the second,
\(P(\mathrm C)=\tfrac12(.1)+\tfrac12(.8)=.45\),
so \(P(\neg\mathrm C)=.55\); conditioning gives \(.2\).
Both evidence events have probability \(1/2\), and all four
prior assignments have positive mass. The contrast therefore
does not exploit conditioning on an impossible event.
\end{proof}
\end{annexitem}

\begin{annexitem}{annex:p16.rat42-covariance}{thm:p16-rat42-covariance}{The weighted identity, including the countable case}
\begin{proof}
Normalization gives
\[
 P(\mathrm H)=\sum_i y_i h_i,
 \quad P(\mathrm F)=\sum_i y_i f_i,
 \quad P(\mathrm H,\mathrm F)=\sum_i y_i h_i f_i.
\]
The last equality uses the independent local \((H,F)\) table
conditional on \(Y\). Expanding the symmetric double sum gives
\[
\begin{aligned}
 \tfrac12\sum_{i,j}y_i y_j(h_i-h_j)(f_i-f_j)
 &=\tfrac12\sum_{i,j}y_i y_j
   (h_if_i+h_jf_j-h_if_j-h_jf_i)\\
 &=\sum_i y_i h_i f_i
       -\left(\sum_i y_i h_i\right)
        \left(\sum_j y_j f_j\right).
\end{aligned}
\]
In the countable case, each summand in each expanded expression
is bounded in absolute value by \(y_i y_j\), whose double sum
is one. Thus all these series converge absolutely and their
rearrangement is legitimate; equivalently finite truncations
converge with tails bounded by the omitted \(Y\)-mass.
Subtract \(P(\mathrm H)\) from
\(P(\mathrm H,\mathrm F)/P(\mathrm F)\) and substitute this
identity. The positive denominator is stipulated in the theorem.
\end{proof}
\end{annexitem}

\begin{annexitem}{annex:p16.rat42-naa}{thm:p16-rat42-naa}{Complete strictness criterion for no alternatives}
\begin{proof}
For \(i>j\), monotonicity gives \(h_i-h_j\leq0\) and
\(f_i-f_j\leq0\); their product is nonnegative. The same is
true when \(j>i\), and diagonal terms vanish. The covariance
numerator is therefore nonnegative. It is strictly positive
exactly when at least one term has positive weight and two
strict differences. Orienting its indices as \(i>j\) gives
\(y_i y_j>0,h_i<h_j,f_i<f_j\). Divide by the positive
\(P(\mathrm F)\) to obtain the theorem. If no such pair exists,
every term is zero, so there is no confirmation change.

For the two-state illustration,
\(P(\mathrm H)=P(\mathrm F)=1/2\), and
\[
 P(\mathrm H,\mathrm F)
       =\tfrac12(4/5)(4/5)+\tfrac12(1/5)(1/5)
       =17/50.
\]
Its posterior is \((17/50)/(1/2)=17/25\), strictly above
the prior. NAA1's uncertainty restriction is respected
by both weights \(1/2\); the proof also shows why that
restriction alone, without the two varying conditional tables,
would not suffice.
\end{proof}
\end{annexitem}

\begin{annexitem}{annex:p16.rat42-nrf}{thm:p16-rat42-nrf}{Complete strictness criterion for no reason for}
\begin{proof}
For \(i>j\), NRF2 gives \(h_i-h_j\geq0\), while NRF3
gives \(f_i-f_j\leq0\). Every off-diagonal covariance term
is nonpositive. Its sum is negative exactly when some
term has positive weight and both differences are strict:
\(y_i y_j>0,h_i>h_j,f_i<f_j\). This holds for both finite
and countable values; a sum of nonpositive absolutely
summable terms is zero exactly when each is zero.
The denominator is positive, so the covariance identity
gives precisely the disconfirmation criterion.

In the numerical illustration both marginals are \(1/2\),
while
\[
 P(\mathrm H,\mathrm F)
       =\tfrac12(1/5)(4/5)+\tfrac12(4/5)(1/5)=4/25.
\]
Hence the posterior is \(8/25<1/2\). In this legitimate
example \(i=1>0=j\) has \(f_1<f_0\) and \(h_1>h_0\),
exactly the stated directions. These are the signs required
by the model's assumptions and its computed change.
\end{proof}
\end{annexitem}

\begin{annexitem}{annex:p16.rat42-expertise-trust}{thm:p16-rat42-expertise-trust}{Separating ignorance from informed deceit}
\begin{proof}
Under a true hypothesis, the honest-expert branch has probability
\(et\) and gives a positive report with certainty; the
dishonest-expert branch gives none. The nonexpert branches
together have probability \(1-e\) and each gives positive
report probability \(1/2\), yielding \(p=et+(1-e)/2\).
Under a false hypothesis, precisely the dishonest expert
branch of weight \(e(1-t)\) gives a positive report;
the nonexpert contribution is unchanged. Thus
\(q=e(1-t)+(1-e)/2\), and subtraction gives
\(p-q=e(2t-1)\). The confirmation lemma identifies its
sign with the posterior change whenever \(0<h<1\)
and the report is positive-probability. These likelihoods
satisfy \(p+q=1\), so the evidence probability is indeed
positive under that interior prior.

At \(e=4/5,t=9/10\), the likelihoods are
\((4/5)(9/10)+1/10=41/50\) and
\((4/5)(1/10)+1/10=9/50\). At \(t=1/10\)
they reverse. The mathematical result depends on the
specified deceit table; it does not assert that every
untrustworthy real source negates the truth deterministically.
\end{proof}
\end{annexitem}
